\chapter{1992年三南卷}
\section{单选题}

\begin{problem}
设函数 \(z=\i^2+\sqrt{3}\i\)，那么 \(\arg z\) 是（\quad）

\choices
    {\(\frac{5}{6}\pi\)}
    {\(\frac{\pi}{3}\)}
    {\(\frac{2}{3}\pi\)}
    {\(-\frac{4}{3}\pi\)}
\end{problem}

\begin{answer}
C
\end{answer}

\begin{solution}
由 \(\i^2=-1\)，得 \(z=-1+\sqrt{3}\i\)，所以 \(z\) 在第二象限，且
\[
\tan\arg z=\frac{\sqrt{3}}{-1}=-\sqrt{3}.
\]
因此主值 \(\arg z=\frac{2}{3}\pi\)，故选 C.
\end{solution}

\begin{problem}
如果等边圆柱（即底面直径与母线相等的圆柱）的体积是 \(16\pi\mathrm{~cm}^3\)，那么它的底面半径等于（\quad）

\choices
    {\(4\sqrt[3]{2}\mathrm{~cm}\)}
    {\(4\mathrm{~cm}\)}
    {\(2\sqrt[3]{2}\mathrm{~cm}\)}
    {\(2\mathrm{~cm}\)}
\end{problem}

\begin{answer}
D
\end{answer}

\begin{solution}
设底面半径为 \(r\)，则等边圆柱的高等于底面直径，即 \(h=2r\). 由圆柱体积公式，
\[
16\pi=\pi r^2h=2\pi r^3,
\]
所以 \(r^3=8\)，从而 \(r=2\mathrm{~cm}\)，故选 D.
\end{solution}

\begin{problem}
\(\frac{\arcsin \frac{\sqrt{3}}{2}-\arccos \left( -\frac{1}{2} \right)}{\arctan \left( -\sqrt{3} \right)}\) 的值等于（\quad）

\choices
    {\(1\)}
    {\(0\)}
    {\(-\frac{2}{5}\)}
    {\(\frac{6}{5}\)}
\end{problem}

\begin{answer}
A
\end{answer}

\begin{solution}
由反三角函数的主值范围，
\[
\arcsin\frac{\sqrt{3}}{2}=\frac{\pi}{3},\qquad
\arccos\left(-\frac12\right)=\frac{2\pi}{3},\qquad
\arctan(-\sqrt{3})=-\frac{\pi}{3}.
\]
因此原式等于 \(\dfrac{-\frac{\pi}{3}}{-\frac{\pi}{3}}=1\)，故选 A.
\end{solution}

\begin{problem}
函数 \(y=\log_{\frac{1}{2}}(1-x)\)（\(x<1\)）的反函数是（\quad）

\choices
    {\(y=1+2^{-x}\)（\(x\in\R\)）}
    {\(y=1-2^{-x}\)（\(x\in\R\)）}
    {\(y=1+2^x\)（\(x\in\R\)）}
    {\(y=1-2^x\)（\(x\in\R\)）}
\end{problem}

\begin{answer}
B
\end{answer}

\begin{solution}
交换原函数中的 \(x\) 与 \(y\)，得 \(x=\log_{\frac12}(1-y)\). 于是
\[
\left(\frac12\right)^x=1-y,
\]
所以反函数为 \(y=1-2^{-x}\)，故选 B.
\end{solution}

\begin{problem}
在长方体 \(ABCD-A_1B_1C_1D_1\) 中，如果 \(AB=BC=a\)，\(AA_1=2a\)，那么点 \(A\) 到直线 \(A_1C\) 的距离等于（\quad）

\choices
    {\(\frac{2\sqrt{6}}{3}a\)}
    {\(\frac{3\sqrt{6}}{2}a\)}
    {\(\frac{2\sqrt{3}}{3}a\)}
    {\(\frac{\sqrt{6}}{3}a\)}
\end{problem}

\begin{answer}
C
\end{answer}

\begin{solution}
取 \(A\) 为原点，令 \(AB\)、\(BC\)、\(AA_1\) 分别沿三个坐标轴正向，则
\[
A_1=(0,0,2a),\qquad C=(a,a,0).
\]
直线 \(A_1C\) 的方向向量可取 \(\bs{v}=(a,a,-2a)\)，而 \(\overrightarrow{A_1A}=(0,0,-2a)\). 点到直线的距离为
\[
d=\frac{\left|\overrightarrow{A_1A}\times\bs{v}\right|}{|\bs{v}|}
=\frac{2\sqrt{2}a^2}{\sqrt{6}a}
=\frac{2\sqrt{3}}{3}a.
\]
故选 C.
\end{solution}

\begin{problem}
函数 \(y=\sin x\cos x+\sqrt{3}\cos^2 x-\frac{\sqrt{3}}{2}\) 的最小正周期等于（\quad）

\choices
    {\(\pi\)}
    {\(2\pi\)}
    {\(\frac{\pi}{4}\)}
    {\(\frac{\pi}{2}\)}
\end{problem}

\begin{answer}
A
\end{answer}

\begin{solution}
利用二倍角公式，原函数化为
\[
y=\frac12\sin2x+\frac{\sqrt{3}}{2}(1+\cos2x)-\frac{\sqrt{3}}2
=\frac12\sin2x+\frac{\sqrt{3}}2\cos2x
=\sin\left(2x+\frac{\pi}{3}\right).
\]
因此其最小正周期为 \(\dfrac{2\pi}{2}=\pi\)，故选 A.
\end{solution}

\begin{problem}
有一个椭圆，它的极坐标方程是（\quad）

\choices
    {\(\rho=\frac{5}{\sqrt{3}-2\cos\theta}\)}
    {\(\rho=\frac{5}{\sqrt{3}-\sqrt{3}\cos\theta}\)}
    {\(\rho=\frac{2-\sqrt{3}\cos\theta}{5}\)}
    {\(\rho=\frac{5}{2-\sqrt{3}\cos\theta}\)}
\end{problem}

\begin{answer}
D
\end{answer}

\begin{solution}
椭圆的极坐标方程可写成 \(\rho=\dfrac{p}{1-e\cos\theta}\)，其中 \(0<e<1\). 选项 D 可化为
\[
\rho=\frac{5/2}{1-\frac{\sqrt{3}}2\cos\theta},
\]
其离心率为 \(e=\frac{\sqrt{3}}2<1\)，所以表示椭圆，故选 D.
\end{solution}

\begin{problem}
不等式 \(\left| \sqrt{x-2}-3 \right|<1\) 的解集是（\quad）

\choices
    {\(\{x\mid 5<x<16\}\)}
    {\(\{x\mid 6<x<18\}\)}
    {\(\{x\mid 7<x<20\}\)}
    {\(\{x\mid 8<x<22\}\)}
\end{problem}

\begin{answer}
B
\end{answer}

\begin{solution}
由根式有意义知 \(x\geqslant2\). 在此条件下，
\[
\left|\sqrt{x-2}-3\right|<1
\Longleftrightarrow 2<\sqrt{x-2}<4.
\]
两边平方得 \(4<x-2<16\)，即 \(6<x<18\)，故选 B.
\end{solution}

\begin{problem}
设等差数列 \(\{a_n\}\) 的公差是 \(d\)，如果它的前 \(n\) 项和 \(S_n=-n^2\)，那么（\quad）

\choices
    {\(a_n=2n-1\)，\(d=-2\)}
    {\(a_n=2n-1\)，\(d=2\)}
    {\(a_n=-2n+1\)，\(d=-2\)}
    {\(a_n=-2n+1\)，\(d=2\)}
\end{problem}

\begin{answer}
C
\end{answer}

\begin{solution}
由等差数列前 \(n\) 项和的关系，
\[
a_n=S_n-S_{n-1}=-n^2+(n-1)^2=-2n+1.
\]
因此 \(d=a_{n+1}-a_n=-2\)，故选 C.
\end{solution}

\begin{problem}
方程 \(\cos 2x=3\cos x+1\) 的解集是（\quad）

\choices
    {\(\left\{x\mid x=2k\pi\pm\frac{2}{3}\pi,\ k\in\Z\right\}\)}
    {\(\left\{x\mid x=k\pi\pm\frac{1}{3}\pi,\ k\in\Z\right\}\)}
    {\(\left\{x\mid x=k\pi\pm\frac{2}{3}\pi,\ k\in\Z\right\}\)}
    {\(\left\{x\mid x=2k\pi\pm\frac{1}{3}\pi,\ k\in\Z\right\}\)}
\end{problem}

\begin{answer}
A
\end{answer}

\begin{solution}
令 \(u=\cos x\)，由 \(\cos2x=2\cos^2x-1\) 得
\[
2u^2-1=3u+1,\qquad 2u^2-3u-2=0.
\]
即 \((2u+1)(u-2)=0\). 因为 \(|u|\leqslant1\)，只能有 \(u=-\frac12\)，所以
\[
x=2k\pi\pm\frac{2}{3}\pi,\qquad k\in\Z.
\]
故选 A.
\end{solution}

\begin{problem}
有一条半径是 \(2\) 的弧，其度数是 \(60^\circ\)，它绕经过弧的中点的直径旋转得到一个球冠，那么这个球冠的面积是（\quad）

\choices
    {\(4(2-\sqrt{3})\pi\)}
    {\(2(2-\sqrt{3}\pi)\)}
    {\(4\sqrt{3}\pi\)}
    {\(2\sqrt{3}\pi\)}
\end{problem}

\begin{answer}
A
\end{answer}

\begin{solution}
球冠所在球面的半径为 \(R=2\). 弧的圆心角为 \(60^\circ\)，绕经过弧中点的直径旋转后，球冠的极角为 \(30^\circ\)，所以球冠高为
\[
h=R-R\cos30^\circ=2-\sqrt{3}.
\]
球冠面积为 \(2\pi Rh=2\pi\cdot2(2-\sqrt{3})=4(2-\sqrt{3})\pi\)，故选 A.
\end{solution}

\begin{problem}
某小组共有 \(10\) 名学生，其中女生 \(3\) 名. 现选举 \(2\) 名代表，至少有 \(1\) 名女生当选的不同的选法共有（\quad）

\choices
    {\(27\) 种}
    {\(48\) 种}
    {\(21\) 种}
    {\(24\) 种}
\end{problem}

\begin{answer}
D
\end{answer}

\begin{solution}
按当选女生人数分类：恰有 \(1\) 名女生时有 \(\mathrm C_3^1\mathrm C_7^1=21\) 种，恰有 \(2\) 名女生时有 \(\mathrm C_3^2=3\) 种. 因此至少有 \(1\) 名女生当选的选法共有 \(21+3=24\) 种，故选 D.
\end{solution}

\begin{problem}
设全集 \(U=\R\)，集合 \(M=\left\{x\mid \sqrt{x^2}>2\right\}\)，\(N=\left\{x\mid \log_x 7>\log_3 7\right\}\)，那么 \(M\cap\overline{N}=\)（\quad）

\choices
    {\(\{x\mid x<-2\}\)}
    {\(\{x\mid x<-2\text{\ 或\ }x\geqslant 3\}\)}
    {\(\{x\mid x\geqslant 3\}\)}
    {\(\{x\mid -2\leqslant x<3\}\)}
\end{problem}

\begin{answer}
B
\end{answer}

\begin{solution}
因为 \(\sqrt{x^2}=|x|\)，所以
\[
M=\{x\mid x<-2\text{ 或 }x>2\}.
\]
对集合 \(N\)，定义域要求 \(x>0\) 且 \(x\ne1\). 当 \(0<x<1\) 时，\(\log_x7<0<\log_37\)；当 \(x>1\) 时，由 \(\log_x7=\dfrac{\ln7}{\ln x}\) 得
\[
\log_x7>\log_37\Longleftrightarrow \ln x<\ln3\Longleftrightarrow 1<x<3.
\]
故 \(N=(1,3)\)，从而 \(\overline N=(-\infty,1]\cup[3,+\infty)\). 因此
\[
M\cap\overline N=(-\infty,-2)\cup[3,+\infty),
\]
故选 B.
\end{solution}

\begin{problem}
设 \(\{a_n\}\) 是由正数组成的等比数列，公比 \(q=2\)，且 \(a_1a_2a_3\cdots a_{30}=2^{30}\)，那么 \(a_3a_6a_9\cdots a_{30}\) 等于（\quad）

\choices
    {\(2^{10}\)}
    {\(2^{20}\)}
    {\(2^{16}\)}
    {\(2^{15}\)}
\end{problem}

\begin{answer}
B
\end{answer}

\begin{solution}
设 \(a_1=A\)，则 \(a_n=A2^{n-1}\). 由已知条件，
\[
A^{30}2^{0+1+\cdots+29}=2^{30},
\]
所以 \(A^{30}=2^{-405}\)，因数列各项为正，得 \(A=2^{-27/2}\). 于是
\[
a_3a_6\cdots a_{30}=A^{10}2^{(3-1)+(6-1)+\cdots+(30-1)}
=2^{-135}2^{155}=2^{20}.
\]
故选 B.
\end{solution}

\begin{problem}
设 \(\triangle ABC\) 不是直角三角形，\(A\) 和 \(B\) 是它的两个内角，那么（\quad）

\choices
    {“\(A<B\)”是“\(\tan A<\tan B\)”的充分条件，但不是必要条件}
    {“\(A<B\)”是“\(\tan A<\tan B\)”的必要条件，但不是充分条件}
    {“\(A<B\)”是“\(\tan A<\tan B\)”的充分必要条件}
    {“\(A<B\)”不是“\(\tan A<\tan B\)”的充分条件，也不是必要条件}
\end{problem}

\begin{answer}
D
\end{answer}

\begin{solution}
若取 \(A=30^\circ\)，\(B=120^\circ\)，则 \(A<B\)，但 \(\tan A=\frac{\sqrt{3}}3>-\sqrt{3}=\tan B\)，所以 \(A<B\) 不是充分条件. 若取 \(A=120^\circ\)，\(B=30^\circ\)，则 \(\tan A=-\sqrt{3}<\frac{\sqrt{3}}3=\tan B\)，但 \(A<B\) 不成立，所以 \(A<B\) 不是必要条件. 故选 D.
\end{solution}

\begin{problem}
对于定义域是 \(\R\) 的任何奇函数 \(f(x)\)，都有（\quad）

\choices
    {\(f(x)-f(-x)>0\)（\(x\in\R\)）}
    {\(f(x)-f(-x)\leqslant 0\)（\(x\in\R\)）}
    {\(f(x)f(-x)\leqslant 0\)（\(x\in\R\)）}
    {\(f(x)f(-x)>0\)（\(x\in\R\)）}
\end{problem}

\begin{answer}
C
\end{answer}

\begin{solution}
因为 \(f\) 是奇函数，所以 \(f(-x)=-f(x)\). 因此
\[
f(x)f(-x)=-f(x)^2\leqslant0,
\]
对任意 \(x\in\R\) 都成立，故选 C.
\end{solution}

\begin{problem}
如果双曲线的两条渐近线的方程是 \(y=\pm\frac{3}{2}x\)，焦点坐标是 \((-\sqrt{26},0)\) 和 \((\sqrt{26},0)\)，那么它的两条准线之间的距离是（\quad）

\choices
    {\(\frac{8}{13}\sqrt{26}\)}
    {\(\frac{4}{13}\sqrt{26}\)}
    {\(\frac{18}{13}\sqrt{26}\)}
    {\(\frac{9}{13}\sqrt{26}\)}
\end{problem}

\begin{answer}
A
\end{answer}

\begin{solution}
设双曲线方程为 \(\dfrac{x^2}{a^2}-\dfrac{y^2}{b^2}=1\). 由渐近线斜率得
\[
\frac{b}{a}=\frac32,\qquad b^2=\frac94a^2.
\]
又因焦点横坐标为 \(\sqrt{26}\)，所以
\[
a^2+b^2=c^2=26,\qquad \frac{13}{4}a^2=26,\qquad a^2=8.
\]
双曲线准线为 \(x=\pm\dfrac{a^2}{c}\)，两准线距离为
\[
\frac{2a^2}{c}=\frac{16}{\sqrt{26}}=\frac{8\sqrt{26}}{13},
\]
故选 A.
\end{solution}

\section{填空题}

\begin{problem}
\(\tan \frac{\pi}{8}=\)\fillinblank{}.
\end{problem}

\begin{answer}
\(\sqrt{2}-1\)
\end{answer}

\begin{solution}
由半角公式，
\[
\tan\frac{\pi}{8}=\sqrt{\frac{1-\cos\frac{\pi}{4}}{1+\cos\frac{\pi}{4}}}
=\sqrt{\frac{1-\frac{\sqrt{2}}2}{1+\frac{\sqrt{2}}2}}
=\sqrt{\left(\sqrt{2}-1\right)^2}=\sqrt{2}-1.
\]
\end{solution}

\begin{problem}
设直线的参数方程是 \(\begin{cases}
x=2+\frac{1}{2}t,\\
y=3+\frac{\sqrt{3}}{2}t,
\end{cases}\) 那么它的斜截式方程是\fillinblank{}.
\end{problem}

\begin{answer}
\(y=\sqrt{3}x+3-2\sqrt{3}\)
\end{answer}

\begin{solution}
由参数方程消去参数 \(t\)：
\[
t=2(x-2),
\]
代入 \(y=3+\frac{\sqrt{3}}2t\)，得
\[
y=3+\frac{\sqrt{3}}2\cdot2(x-2)=\sqrt{3}x+3-2\sqrt{3}.
\]
所以斜截式方程为 \(y=\sqrt{3}x+3-2\sqrt{3}\).
\end{solution}

\begin{problem}
如果三角形的顶点分别是 \(O(0,0)\)，\(A(0,15)\)，\(B(-8,0)\)，那么它的内切圆方程是\fillinblank{}.
\end{problem}

\begin{answer}
\((x+3)^2+(y-3)^2=9\)
\end{answer}

\begin{solution}
三角形三边长为
\[
OA=15,\qquad OB=8,\qquad AB=\sqrt{8^2+15^2}=17.
\]
它是直角三角形，直角顶点为 \(O\). 设内心为 \(I\)，按内心坐标公式，
\[
I=\frac{17O+8A+15B}{17+8+15}=(-3,3).
\]
内切圆半径等于内心到坐标轴的距离，即 \(r=3\). 因此内切圆方程为
\[
(x+3)^2+(y-3)^2=3^2,
\]
即 \((x+3)^2+(y-3)^2=9\).
\end{solution}

\begin{problem}
\(\displaystyle\lim_{n\to\infty}\left[ \frac{1}{1\cdot 4}+\frac{1}{4\cdot 7}+\frac{1}{7\cdot 10}+\cdots+\frac{1}{(3n-2)(3n+1)} \right]=\)\fillinblank{}.
\end{problem}

\begin{answer}
\(\frac13\)
\end{answer}

\begin{solution}
对任意正整数 \(k\)，
\[
\frac{1}{(3k-2)(3k+1)}=\frac13\left(\frac{1}{3k-2}-\frac{1}{3k+1}\right).
\]
所以前 \(n\) 项和为
\[
\frac13\left(1-\frac{1}{3n+1}\right).
\]
令 \(n\to\infty\)，得极限为 \(\frac13\).
\end{solution}

\begin{problem}
\(91^{92}\) 除以 \(100\) 的余数是\fillinblank{}.
\end{problem}

\begin{answer}
81
\end{answer}

\begin{solution}
因为 \(91=1-10\)，由二项式定理，
\[
91^{92}=(1-10)^{92}\equiv1-92\times10\pmod {100},
\]
其余各项均含有因子 \(10^2=100\). 因而
\[
91^{92}\equiv1-920\equiv81\pmod {100},
\]
余数为 \(81\).
\end{solution}

\begin{problem}
已知三棱锥 \(A-BCD\) 的体积是 \(V\)，棱 \(BC\) 的长是 \(a\)，面 \(ABC\) 和面 \(DBC\) 的面积分别是 \(S_1\) 和 \(S_2\). 设面 \(ABC\) 和面 \(DBC\) 所成的二面角是 \(\alpha\)，那么 \(\sin \alpha=\)\fillinblank{}.
\end{problem}

\begin{answer}
\(\frac{3aV}{2S_1S_2}\)
\end{answer}

\begin{solution}
设面 \(ABC\) 内垂直于 \(BC\) 的高为 \(h_1\)，面 \(DBC\) 内垂直于 \(BC\) 的高为 \(h_2\). 由三角形面积公式，
\[
h_1=\frac{2S_1}{a},\qquad h_2=\frac{2S_2}{a}.
\]
沿棱 \(BC\) 作垂直截面，截面内两条高的夹角为 \(\alpha\)，故三棱锥体积为
\[
V=\frac16ah_1h_2\sin\alpha.
\]
代入 \(h_1\)、\(h_2\)，得
\[
V=\frac{2S_1S_2}{3a}\sin\alpha,
\]
所以 \(\sin\alpha=\frac{3aV}{2S_1S_2}\).
\end{solution}

\section{解答题}

\begin{problem}
已知关于 \(x\) 的方程 \(2a^{2x-2}-7a^{x-1}+3=0\) 有一个根是 \(2\)，求 \(a\) 的值和方程其余的根.
\end{problem}

\begin{solution}
指数式有意义且为指数函数的底数，须有 \(a>0\) 且 \(a\ne1\). 令 \(t=a^{x-1}\)，原方程化为
\[
2t^2-7t+3=0,
\]
即
\[
(2t-1)(t-3)=0.
\]
因此 \(t=\frac12\) 或 \(t=3\). 已知 \(x=2\) 是方程的一个根，代入得 \(t=a\)，所以 \(a=\frac12\) 或 \(a=3\).

当 \(a=\frac12\) 时，\(t=\frac12\) 给出已知根 \(x=2\)，另一个根满足
\[
\left(\frac12\right)^{x-1}=3,
\]
故 \(x-1=\log_{\frac12}3\)，即 \(x=1+\log_{\frac12}3\).

当 \(a=3\) 时，\(t=3\) 给出已知根 \(x=2\)，另一个根满足
\[
3^{x-1}=\frac12,
\]
故 \(x-1=\log_3\frac12=-\log_3 2\)，即 \(x=1-\log_3 2\).
\end{solution}

\begin{problem}
已知平面 \(\alpha\) 和不在这个平面内的直线 \(a\) 都垂直于平面 \(\beta\). 求证： \(a\parallel \alpha\).
\end{problem}

\begin{solution}
设 \(l=\alpha\cap\beta\). 在平面 \(\alpha\) 内过任意一点作直线 \(m\)，使 \(m\perp l\). 由于 \(\alpha\perp\beta\)，二面角的平面角为直角，所以 \(m\perp\beta\). 题设又给出 \(a\perp\beta\)，因此两条同垂直于平面的直线平行，即 \(a\parallel m\).

因为 \(m\subset\alpha\)，所以 \(a\) 的方向平行于平面 \(\alpha\). 若 \(a\) 与 \(\alpha\) 相交，则过交点且与 \(m\) 平行的直线应在平面 \(\alpha\) 内，从而 \(a\subset\alpha\)，这与直线 \(a\) 不在平面 \(\alpha\) 内矛盾. 因此 \(a\) 与 \(\alpha\) 不相交，故 \(a\parallel\alpha\).
\end{solution}

\begin{problem}
证明不等式： \(1+\frac{1}{\sqrt{2}}+\frac{1}{\sqrt{3}}+\cdots+\frac{1}{\sqrt{n}}<2\sqrt{n}\)（\(n\in \N^*\)）.
\end{problem}

\begin{solution}
对任意正整数 \(k\)，有
\[
2\left(\sqrt{k}-\sqrt{k-1}\right)=\frac{2}{\sqrt{k}+\sqrt{k-1}}>\frac{1}{\sqrt{k}},
\]
因为 \(2\sqrt{k}>\sqrt{k}+\sqrt{k-1}\). 因此逐项相加，得到
\[
1+\frac{1}{\sqrt{2}}+\cdots+\frac{1}{\sqrt{n}}
<2\sum_{k=1}^{n}\left(\sqrt{k}-\sqrt{k-1}\right)=2\sqrt{n}.
\]
不等式得证.
\end{solution}

\begin{problem}
设抛物线经过两点 \((-1,6)\) 和 \((-1,-2)\)，对称轴与 \(x\) 轴平行，开口向右；直线 \(y=2x+7\) 被抛物线截得的线段的长是 \(4\sqrt{10}\)，求抛物线的方程.
\end{problem}

\begin{solution}
设抛物线方程为
\[
(y-k)^2=4p(x-h),\qquad p>0.
\]
由于点 \((-1,6)\) 和 \((-1,-2)\) 的纵坐标关于对称轴对称，得
\[
k=\frac{6+(-2)}2=2.
\]
代入点 \((-1,6)\)，得
\[
16=4p(-1-h),\qquad p(-1-h)=4.
\]

直线与抛物线的交点横坐标是方程
\[
(2x+7-2)^2=4p(x-h)
\]
的两个根，即
\[
4x^2+(20-4p)x+25+4ph=0.
\]
设两根为 \(x_1\)、\(x_2\). 直线斜率为 \(2\)，所以截得线段长为
\[
\sqrt{1+2^2}|x_1-x_2|=\sqrt{5}|x_1-x_2|=4\sqrt{10},
\]
从而 \(|x_1-x_2|=4\sqrt{2}\). 由二次方程根的差公式，
\[
|x_1-x_2|^2=p(p-10-4h).
\]
故
\[
p(p-10-4h)=32.
\]
由 \(h=-1-\frac4p\)，得
\[
p^2-6p+16=32,\qquad p^2-6p-16=0.
\]
因 \(p>0\)，所以 \(p=8\)，进而 \(h=-\frac32\). 因此抛物线方程为
\[
(y-2)^2=32\left(x+\frac32\right),
\]
即 \((y-2)^2=32x+48\).
\end{solution}

\begin{problem}
求同时满足下列两个条件的所有复数 \(z\)：
\begin{circlelist}
\item \(z+\frac{10}{z}\) 是实数，且 \(1<z+\frac{10}{z}\leqslant 6\)；
\item \(z\) 的实部和虚部都是整数.
\end{circlelist}
\end{problem}

\begin{solution}
设 \(z=x+y\i\)，其中 \(x,y\in\Z\)，且 \(z\ne0\). 因为
\[
z+\frac{10}{z}=x+y\i+\frac{10(x-y\i)}{x^2+y^2},
\]
其虚部为 \(y\left(1-\frac{10}{x^2+y^2}\right)\). 该式为实数，所以分两种情况讨论.

当 \(y=0\) 时，\(x\ne0\)，且题设不等式要求
\[
1<x+\frac{10}{x}\leqslant6.
\]
若 \(x<0\)，左端不可能大于 \(1\)；若 \(x>0\)，则
\[
x+\frac{10}{x}\geqslant2\sqrt{10}>6,
\]
也不满足条件. 因此此时无解.

当 \(x^2+y^2=10\) 时，由 \(x,y\in\Z\)，有
\[
(x,y)=(\pm1,\pm3)\quad\text{或}\quad(\pm3,\pm1).
\]
此时
\[
z+\frac{10}{z}=x+y\i+\frac{10(x-y\i)}{10}=2x.
\]
条件 \(1<2x\leqslant6\) 要求 \(x=1\) 或 \(x=3\). 所以所有复数为
\[
z=1+3\i,\quad 1-3\i,\quad 3+\i,\quad 3-\i.
\]
\end{solution}
